%% 第二章

\chapter{相关理论与关键技术基础}
\chaptermark{相关理论与关键技术基础}
\section{三维重建与新视角合成基础}
三维重建的目标是根据二维图像恢复场景中的几何结构、相机位置和外观信息。对于静态深空场景，如果能够从不同位置采集多张图像，则不同图像之间存在由相机投影关系决定的几何约束。通过建立图像像素、相机参数和三维点之间的对应关系，可以估计相机运动并恢复场景结构。

为了将恢复出的空间信息用于可视化、分析和后续渲染，还需要选择合适的三维场景表示形式。三维重建结果可以表示为点云、网格、体素、隐式场或显式高斯等形式。点云便于描述稀疏空间分布，网格适合连续表面表达，体素便于表示空间占用关系，隐式场具有连续函数表达能力，显式高斯则能够在显式参数化的基础上支持高效可微渲染。不同表示形式在表达能力、计算代价和适用场景上各有侧重，但都需要通过相机成像模型建立三维空间与二维图像之间的投影关系。

在针孔相机模型中，三维点到二维图像平面的投影关系可表示为：
\begin{equation}
    s\left[\begin{smallmatrix}u\\v\\1\end{smallmatrix}\right]=K\left[\,R\mid t\,\right]\left[\begin{smallmatrix}X\\Y\\Z\\1\end{smallmatrix}\right]
\end{equation}
其中，$K$为相机内参矩阵，$R$和$t$分别表示相机旋转和平移，$(X,Y,Z)$为三维点坐标，$(u,v)$为其在图像中的投影坐标。

新视角合成是在已有图像观测和场景表示的基础上，从训练图像中未直接观测到的观察位置生成对应图像的任务。该任务通常以多视角图像和相机参数为输入，通过恢复或学习可渲染的三维场景表示，在新的相机视角下生成图像。与三维重建侧重空间结构恢复不同，新视角合成更强调目标视角下图像内容的生成质量，通常通过渲染图像与真实图像之间的差异进行评价。

\section{SfM与COLMAP相机位姿估计}
SfM即运动恢复结构，是一种根据多张图像同时估计相机运动和场景三维结构的方法。其基本流程包括特征点提取、特征匹配、几何验证、相机位姿估计、三角化和全局优化。COLMAP将这一流程工程化，在三维重建和新视角合成任务中常用于提供相机位姿和初始稀疏点云\upcite{ref2,ref3}。图\ref{fig:2-colmap-rome}给出了COLMAP在无序图像集合上恢复相机位姿与稀疏点云的典型示例。

\begin{figure}[H]
    \centering
    \includegraphics[width=0.82\textwidth]{colmap重建的罗马城.jpg}
    \caption{COLMAP重建的罗马城示例}\label{fig:2-colmap-rome}
\end{figure}

在SfM中，重投影误差是衡量相机位姿和三维点估计精度的重要依据，其基本形式为：
\begin{equation}
    E_{\mathrm{reproj}}=\sum_j\sum_{i\in V_j}\left\|x_{ij}-\pi\left(K_j(R_jX_i+t_j)\right)\right\|_2^2
\end{equation}
Bundle Adjustment通过联合优化相机参数和三维点坐标，使整体重投影误差最小：
\begin{equation}
    \min_{\{R_j,t_j,K_j\},\{X_i\}}\sum_j\sum_{i\in V_j}\rho\left(\left\|x_{ij}-\pi\left(K_j(R_jX_i+t_j)\right)\right\|_2^2\right)
\end{equation}
需要注意的是，基于BA优化的SfM主要恢复相机位姿和场景稀疏点云，其结果并不直接包含完整的连续表面信息。为了获得更高质量的场景表示和新视角合成效果，通常需要在位姿估计和初始点云的基础上，进一步结合MVS等方法进行稠密重建，或利用3D Gaussian Splatting等显式可微渲染技术构建可渲染的三维场景表示。

\section{3D Gaussian Splatting基本原理}
3D Gaussian Splatting是一种基于显式三维高斯基元的场景表示与可微渲染方法。与NeRF等隐式神经表示不同，3DGS不需要沿相机光线进行大量采样，而是直接使用一组三维高斯基元表示场景中的几何结构和外观信息。每个高斯基元通常包含三维中心位置、尺度、旋转、不透明度以及颜色特征等属性。通过对这些属性进行优化，模型可以在保持较高渲染效率的同时实现高质量的新视角合成。

对于第i个三维高斯，其空间中心记为$\mu_i$，协方差矩阵记为$\Sigma_i$。给定三维空间中的任意位置x，该高斯对该位置的影响可表示为：
\begin{equation}
    G_i(x)=\exp\left(-\frac{1}{2}(x-\mu_i)^T\Sigma_i^{-1}(x-\mu_i)\right)
\end{equation}
其中，$\mu_i$决定高斯在三维空间中的位置，$\Sigma_i$决定高斯的空间形态和覆盖范围。由于协方差矩阵需要保持半正定，3DGS通常不直接优化完整协方差矩阵，而是将其分解为旋转矩阵和尺度矩阵的形式：
\begin{equation}
    \Sigma_i=R_iS_iS_i^TR_i^T
\end{equation}
其中，$R_i$表示第i个高斯的旋转矩阵，$S_i$表示尺度矩阵。尺度参数控制高斯在三个主方向上的伸展程度，旋转参数控制这些主方向在三维空间中的朝向。通过这种分解方式，可以保证协方差矩阵具有合理的几何性质，同时使高斯基元能够表达各向异性的空间结构。

在渲染过程中，三维高斯需要根据当前相机参数投影到二维图像平面。设当前视角下的观察变换矩阵为$W$，投影变换在局部仿射近似下的雅可比矩阵为$J$，则三维协方差矩阵经过投影后可得到二维屏幕空间协方差：
\begin{equation}
    \Sigma'_i=JW\Sigma_iW^TJ^T
\end{equation}
该二维协方差决定了高斯投影到图像平面后的椭圆形覆盖区域。相比将三维点简单投影为单个像素点，3DGS将每个高斯投影为具有一定空间范围的二维椭圆分布，因此能够以splatting的方式高效完成图像渲染。屏幕空间中的高斯会根据深度顺序进行排序，并按照透明度进行颜色累积。

每个高斯的颜色通常由球谐函数表示，从而建模一定程度的视角相关外观变化。对于观察方向$d$，第$i$个高斯的颜色可写为：
\begin{equation}
    c_i(d)=\sum_{l=0}^{L}\sum_{m=-l}^{l}c_{i,lm}Y_{lm}(d)
\end{equation}
其中，$Y_{lm}(d)$表示球谐基函数，$c_{i,lm}$表示对应的球谐系数，$L$为球谐函数的阶数。球谐表示使高斯颜色能够随观察方向发生变化，因此可以在一定程度上表示非朗伯表面、视角相关亮度变化和局部高光等外观现象。

在得到屏幕空间高斯分布和视角相关颜色后，3DGS采用基于透明度的前向alpha blending计算像素颜色。对于像素$p$，设覆盖该像素的高斯集合为$\mathcal{N}(p)$，并按照深度从近到远排序，则该像素的颜色可表示为：
\begin{equation}
    C(p)=\sum_{i\in\mathcal{N}(p)}T_i\alpha_i c_i
\end{equation}
其中，透射率$T_i$表示位于其前方的高斯累积后剩余的透射率，可写为：
\begin{equation}
    T_i=\prod_{j<i}(1-\alpha_j)
\end{equation}
式中，$\alpha_i$表示第$i$个高斯在该像素处的不透明度贡献，$c_i$表示该高斯在当前观察方向下的颜色。该公式表明，靠近相机且不透明度较高的高斯会对最终像素颜色产生更大影响，而后方高斯的贡献会被前方高斯逐渐遮挡。

类似地，3DGS也可以通过相同的透明度累积方式渲染深度图。若$d_i$表示第$i$个高斯在当前视角下的深度，则像素$p$处的渲染深度可表示为：
\begin{equation}
    D(p)=\sum_{i\in\mathcal{N}(p)}T_i\alpha_i d_i
\end{equation}
深度渲染结果能够反映三维高斯场景在当前视角下的空间层次关系。需要注意的是，该深度并不一定等同于真实场景表面深度，而是由多个高斯按照透明度加权累积得到的渲染深度。

综上，每个三维高斯基元可以表示为一组可优化参数：
\begin{equation}
    \theta_i=\{\mu_i,s_i,r_i,\alpha_i,sh_i\}
\end{equation}
其中，$\mu_i$表示高斯中心位置，$s_i$表示尺度参数，$r_i$表示旋转参数，$\alpha_i$表示不透明度，$sh_i$表示球谐颜色系数。在训练过程中，3DGS通过可微光栅化器将三维高斯渲染为图像，并将渲染图像与真实图像进行比较，利用反向传播优化上述高斯参数。

基础图像重建损失通常由$\mathcal{L}_1$像素误差和结构相似性损失共同组成：
\begin{equation}
    \mathcal{L}_{\mathrm{base}}=(1-\lambda_{\mathrm{ssim}})\mathcal{L}_1+\lambda_{\mathrm{ssim}}\left(1-\mathrm{SSIM}(\hat I,I)\right)
\end{equation}
其中，$\hat I$表示渲染图像，$I$表示真实图像，$\lambda_{\mathrm{ssim}}$为结构相似性损失权重。$\mathcal{L}_1$损失关注像素级颜色误差，SSIM损失关注局部结构相似性，二者结合能够同时约束颜色重建和图像结构保持。

3DGS的优势在于显式表示直观、训练速度较快、渲染效率较高，并能够通过可微splatting实现高质量实时新视角合成。与基于隐式神经场的方法相比，3DGS直接维护一组三维高斯基元，其空间位置、尺度、旋转、不透明度和颜色特征均具有明确的参数含义，便于进行可视化分析和后续渲染操作。同时，基于光栅化的前向渲染过程避免了沿每条相机光线进行大量采样，使其在训练效率和实时显示方面具有明显优势。

\section{深空场景数据集介绍}
深空场景数据通常来源于真实探测任务归档、轨道遥感观测、巡视器近地表成像、人工语义标注以及仿真合成流程。与常规室内场景或物体级数据相比，深空场景数据具有成像距离跨度大、纹理信息弱、光照变化明显、视角覆盖稀疏和任务约束强等特点。因此，这类数据往往不能只依赖图像本身，还需要结合相机参数、时间序列关系、语义标注、地形高程或三维几何信息，才能支撑三维重建、新视角合成、导航感知和任务仿真等研究。

根据数据来源、组织形式和任务目标的不同，深空场景数据大致可分为以下五类。

第一类是真实任务归档数据，代表性数据源包括NASA Planetary Data System (PDS)和欧洲航天局Planetary Science Archive (PSA)\upcite{ref33,ref35}。这类数据主要保存探测任务返回的科学与工程数据，覆盖轨道器、着陆器、巡视器和其他行星探测任务，具有真实性强、任务覆盖广和长期可追溯等特点。但真实任务归档通常面向科学数据保存和检索，而不是直接面向计算机视觉任务构建训练集、测试集和统一相机轨迹。因此，在用于三维重建或新视角合成前，通常需要额外完成数据筛选、图像预处理、相机参数整理、视角关系构建和序列划分等工作。

\begin{figure}[H]
    \centering
    \includegraphics[width=0.50\textwidth]{PDS数据图片.jpg}
    \caption{PDS真实任务归档数据示例}\label{fig:2-pds-data}
\end{figure}

第二类是面向多视角重建、三维一致性建模或新视角合成的数据集。M3arsSynth是这类数据的代表之一，其以PDS中的真实火星导航图像为基础，通过三维重建和渲染流程生成多视角视频序列，并提供相机轨迹及相关多模态信息\upcite{ref10}。与仅包含单张图像或语义标签的数据集相比，这类数据具有更明确的视角序列关系和相机运动信息，能够支持训练视角与测试视角的划分，因此更适合用于分析多视角几何约束、连续视角渲染和新视角合成效果。

\begin{figure}[H]
    \centering
    \includegraphics[width=0.82\textwidth]{M3arsSynth视频序列.png}
    \caption{M3arsSynth视频序列数据示例}\label{fig:2-m3arssynth}
\end{figure}

第三类是轨道遥感与地形高程数据。火星、月球和小行星等天体的轨道遥感影像通常具有较大覆盖范围，可用于地表形貌分析、数字高程模型构建和地形理解。例如，HiRISE等火星轨道相机能够提供高分辨率地表影像\upcite{ref36}，Ames Stereo Pipeline (ASP)等立体视觉处理工具可用于从轨道或巡视器立体影像中生成数字地形模型、正射影像和三维模型\upcite{ref37}。这类数据更侧重宏观地形、地貌结构和高程变化，适合用于地形建模、地貌分析和遥感解释，但其成像视角、分辨率和任务组织形式与近地表多视角新视角合成任务仍存在差异。

\begin{figure}[H]
    \centering
    \includegraphics[width=0.82\textwidth]{HIRISE数字地形.jpg}
    \caption{HiRISE数字地形数据示例}\label{fig:2-hirise-terrain}
\end{figure}

第四类是语义标注与场景理解数据。AI4Mars是这类数据的代表，它由Curiosity、Opportunity和Spirit等火星车图像构成，包含大规模像素级地形标注，可用于训练和验证火星场景中的地形类别识别模型\upcite{ref34}。这类数据的优势在于提供了丰富的语义监督，有助于研究障碍物识别、可通行区域分析和任务辅助感知；但其主要目标不是三维重建或新视角合成，通常不提供连续多视角相机轨迹和可直接用于渲染评价的测试视角。因此，这类数据更适合作为深空场景理解和语义先验研究的数据基础。

\begin{figure}[H]
    \centering
    \includegraphics[width=0.75\textwidth]{AImars语义标注数据集.jpg}
    \caption{AI4Mars语义标注数据集示例}\label{fig:2-ai4mars}
\end{figure}

第五类是仿真合成与任务模拟数据。由于真实深空探测数据获取成本高、视角覆盖有限且标注困难，研究中常借助仿真平台、物理渲染器或程序化地形生成方法构造模拟月面、火星地表或小天体表面场景。例如，已有工作通过真实高程模型、撞击坑模型和分形地形构建合成行星地表，用于测试行星着陆与视觉导航系统\upcite{ref38}。这类数据可以较方便地控制相机轨迹、光照方向、地形几何和传感器参数，并能够生成深度、法向、语义标签等多种监督信息，因此适合用于算法预训练、任务演示、导航仿真和极端条件测试。然而，合成数据与真实探测图像之间通常存在外观分布和成像噪声差异，在用于真实深空场景建模时需要关注域差异问题。

\begin{figure}[H]
    \centering
    \includegraphics[width=0.50\textwidth]{simulation.jpg}
    \caption{仿真合成数据示例}\label{fig:2-simulation-data}
\end{figure}

表\ref{tab:2-dataset}进一步总结了上述数据类型的代表数据、主要信息和适用任务。
\begin{table}[H]
    \fontsize{10.5pt}{13.125pt}\selectfont
    \caption{典型深空场景数据类型及其任务侧重点}
    \centering\label{tab:2-dataset}
    \begin{adjustbox}{max width=\textwidth}
    \begin{tabular}{cccc}
        \toprule
        数据类型 & 代表数据或工具 & 主要信息 & 适用任务 \\
        \midrule
        真实任务归档数据 & PDS、PSA & 探测影像、任务元数据 & 真实场景分析、数据筛选、视觉建模基础 \\
        多视角重建与视角合成数据 & M3arsSynth & 多视角图像、相机轨迹、视频序列 & 三维重建、新视角合成、可控视角生成 \\
        轨道遥感与地形高程数据 & HiRISE、DEM数据、ASP & 轨道影像、高程模型、正射影像 & 地形建模、地貌分析、遥感解释 \\
        语义标注与场景理解数据 & AI4Mars & 巡视器图像、像素级语义标签 & 地形分类、可通行区域分析、导航感知 \\
        仿真合成与任务模拟数据 & 合成地形、仿真平台 & 合成图像、深度、法向、语义标签 & 算法预训练、任务模拟、极端条件测试 \\
        \bottomrule
    \end{tabular}
    \end{adjustbox}
\end{table}
总体来看，真实任务归档数据具有较高真实性，多视角重建与视角合成数据更适合三维场景表示和新视角合成研究，轨道遥感与高程数据更适合宏观地形建模和地貌分析，语义标注数据主要服务于场景理解与导航感知，仿真合成数据则便于控制变量和生成多模态监督。不同类型数据在数据来源、相机信息、标注形式、视角组织和适用任务上各有侧重。对于本文关注的三维重建与新视角合成任务而言，多视角图像序列、相机轨迹和可用于渲染评价的目标视角尤为重要，因此具备多视角组织形式的数据更适合支撑后续实验。

\section{评价指标}
新视角合成质量通常通过渲染图像与真实图像之间的差异进行评价。本文采用PSNR、SSIM和LPIPS作为主要评价指标，从像素误差、结构相似性和感知质量三个角度衡量不同方法的新视角合成效果。

PSNR根据均方误差计算，常用于评价渲染图像与真实图像之间的像素级重建误差：
\begin{equation}
    \mathrm{PSNR}=10\log_{10}\frac{\mathrm{MAX}_I^2}{\mathrm{MSE}},\quad \mathrm{MSE}=\frac{1}{HW}\sum_{u,v}\left(\hat I(u,v)-I(u,v)\right)^2
\end{equation}
其中，$\mathrm{MAX}_I$表示图像像素最大取值，$\mathrm{MSE}$表示渲染图像与真实图像之间的均方误差。PSNR数值越高，说明像素级重建误差越小。

SSIM从亮度、对比度和结构三个方面度量图像相似性\upcite{ref28}：
\begin{equation}
    \mathrm{SSIM}(x,y)=\frac{(2\mu_x\mu_y+C_1)(2\sigma_{xy}+C_2)}{(\mu_x^2+\mu_y^2+C_1)(\sigma_x^2+\sigma_y^2+C_2)}
\end{equation}
其中，$\mu_x$和$\mu_y$表示图像局部均值，$\sigma_x$和$\sigma_y$表示局部方差，$\sigma_{xy}$表示协方差。SSIM数值越高，说明渲染图像与真实图像在局部结构上越接近。

LPIPS基于深度特征空间中的距离衡量感知差异\upcite{ref29}：
\begin{equation}
    \mathrm{LPIPS}(x,y)=\sum_l\frac{1}{H_lW_l}\sum_{h,w}\left\|w_l\odot\left(\hat y_{hw}^{\,l}-\hat x_{hw}^{\,l}\right)\right\|_2^2
\end{equation}
其中，$\hat x_{hw}^l$和$\hat y_{hw}^l$表示图像在第$l$层深度特征空间中的归一化特征，$w_l$表示该层特征权重。LPIPS数值越低，说明两幅图像在感知特征上越接近。

除图像质量指标外，本文还使用重投影光度误差和重投影深度误差分析三维场景表示的跨视角一致性。评价时首先将渲染深度与辅助深度进行场景级尺度和偏移对齐，得到用于几何投影的深度图$\tilde D$；随后以每个源视角$s$为起点，选择相机中心距离最近的3个目标视角$t$，构成视角对集合$\mathcal{P}$，并在源视角有效像素中最多采样4096个像素，形成样本集合$\Omega_{s,t}$。对于源视角像素$p=(u,v)$，其三维点可由对齐后的源视角深度反投影得到，并投影到目标视角：
\begin{equation}
    X_s(p)=\Pi_s^{-1}\left(p,\tilde D_s(p)\right),\quad
    p_{s\rightarrow t}=\Pi_t\left(X_s(p)\right),\quad
    z_{s\rightarrow t}(p)=Z_t\left(X_s(p)\right)
\end{equation}
其中，$\Pi_s^{-1}(\cdot)$表示由源视角像素坐标和深度反投影到世界坐标的变换，$\Pi_t(\cdot)$表示由世界坐标投影到目标视角像素坐标的变换，二者均包含对应视角的相机内参和外参；$Z_t(\cdot)$表示三维点在目标视角坐标系下的深度，$z_{s\rightarrow t}(p)$为源视角样本投影到目标视角后的深度。有效样本需要同时满足源视角深度有效、投影深度为正、投影位置位于目标图像范围内，并且目标视角投影位置处的深度有效。

在有效样本集合$\Omega_{s,t}$上，重投影光度误差定义为源视角颜色与目标视角颜色之间的平均绝对差，其中目标视角颜色在非整数投影位置采用双线性采样：
\begin{equation}
    E_{\mathrm{photo}}=\frac{1}{N}\sum_{(s,t)\in\mathcal{P}}\sum_{p\in\Omega_{s,t}}
    \frac{1}{3}\left\|I_s(p)-I_t(p_{s\rightarrow t})\right\|_1
\end{equation}
其中，$N=\sum_{(s,t)\in\mathcal{P}}|\Omega_{s,t}|$为所有有效跨视角样本数量。该指标数值越低，说明当前场景表示在跨视角投影后的颜色对应关系越一致。

重投影深度误差采用目标视角采样深度与投影深度之间的相对误差，其中目标视角深度同样在投影位置进行双线性采样：
\begin{equation}
    E_{\mathrm{depth}}=\frac{1}{N}\sum_{(s,t)\in\mathcal{P}}\sum_{p\in\Omega_{s,t}}
    \frac{\left|\tilde D_t(p_{s\rightarrow t})-z_{s\rightarrow t}(p)\right|}{\max\left(\left|\tilde D_t(p_{s\rightarrow t})\right|,\epsilon\right)}
\end{equation}
其中，$\tilde D_t(p_{s\rightarrow t})$为目标视角对齐深度在投影位置的采样值，$\epsilon$为避免分母过小的常数。该指标数值越低，表示跨视角投影后的深度对应关系越稳定。由于这里使用的是渲染深度与辅助深度对齐后的相对深度关系，因此该指标主要用于辅助分析空间层次稳定性。

\section{本章小结}
本章介绍了三维重建与新视角合成的基本概念，说明了SfM与COLMAP在相机位姿估计和初始点云生成中的作用，重点分析了3D Gaussian Splatting的场景表示、渲染和优化流程。同时，本章概述了PDS、M3arsSynth和AI4Mars等深空场景相关数据集与数据源，并介绍了后续实验采用的PSNR、SSIM、LPIPS、重投影光度误差和重投影深度误差等评价指标。
